\section{The original entrance and the complete local return}
\label{sec:local-results}

The local estimates have three distinct roles.
Proposition~\ref{loc:entrance} derives the third-touchdown state and
its entrance constant from the prescribed physical data.
Proposition~\ref{loc:return} separates the regularity of the exact
reference return, the compliant value error, and the conditions under
which its output can be used as a new input. The proofs of these
estimates occupy Appendices~\ref{sec:entry}--\ref{sec:detuned-return}.

Section~\ref{sec:graph} uses these inputs to compare the complete
return states through the required event indices.
Section~\ref{sec:physical-paths} then reconstructs the ordered paths
and realizes contact prefixes by varying compliance at a fixed
actual contact index. Finally,
Section~\ref{sec:contact-selection} compares the residual entrance
phase with the contact stopping window. This order keeps the local
estimates, the long return comparison and the selection argument
separate, while preserving their common physical trajectory.

% MIXED: actual event labels and proved asymptotic scales; schematic spacing.
% Arrows mean actual event passages (top) or discrete returns (bottom).
\begin{figure}[!htbp]
\centering
\begin{tikzpicture}[
 x=1cm,y=1cm,font=\small,
 every node/.style={inner sep=2pt},
 stage/.style={draw=black!65,rounded corners=2pt,align=center,
               text width=2.35cm,minimum height=1.40cm},
 arrow/.style={-{Latex[length=1.6mm,width=1.1mm]},line width=.5pt},
 chain/.style={draw=black!65,align=center,text width=4.3cm,
               minimum height=.75cm},
 guide/.style={draw=black!40,densely dotted,line width=.5pt}
]
 \node[anchor=west,font=\small\bfseries] at (0,7.05)
   {(a) Actual events from the physical initial state};
 \node[stage] (initial) at (1.25,5.45)
   {original\\preparation};
 \node[stage] (release) at (4.25,5.45)
   {critical release\\$v_r/\varepsilon\to q_{\rm in}>0$};
 \node[stage] (tdone) at (7.25,5.45)
   {signed $\mathrm{TD}_1$\\and contact\\$p_1/\sqrt\varepsilon\to\pi$};
 \node[stage] (tdtwo) at (10.25,5.45)
   {complete $\mathrm{TD}_2$\\$\delta_2\sim C_R\varepsilon^{3/8}$};
 \node[stage] (tdthree) at (13.25,5.45)
   {admitted $\mathrm{TD}_3$\\$\delta_3\asymp\varepsilon^{9/32}$};
 \draw[arrow] (initial.east) -- (release.west);
 \draw[arrow] (release.east) -- (tdone.west);
 \draw[arrow] (tdone.east) -- (tdtwo.west);
 \draw[arrow] (tdtwo.east) -- (tdthree.west);
 \node at (5.75,6.55) {first flight};
 \node at (8.75,6.55) {second flight};
 \node at (11.75,6.55) {one return};
 \node[anchor=west] at (0,4.40) {$\pi$ may have either sign or vanish.};

 \node[anchor=west,font=\small\bfseries] at (0,3.65)
   {(b) A comparison started at the same actual base};
 \node[chain] (actualstart) at (2.35,2.82)
   {actual $(\delta_3,\xi_3^r)$};
 \node[chain] (graphstart) at (2.35,1.65)
   {graph $(\delta_3,\Gamma_\lambda(\delta_3))$};
 \node[chain] (actualend) at (12.05,2.82)
   {actual $(\delta_j,\xi_j^r)$};
 \node[chain] (graphend) at (12.05,1.65)
   {graph $(\bar\delta_j,\Gamma_\lambda(\bar\delta_j))$};
 \draw[arrow] (actualstart.east) -- (6.15,2.82);
 \node at (7.20,2.82) {$\cdots$};
 \draw[arrow] (8.25,2.82) -- (actualend.west);
 \draw[arrow,densely dashed] (graphstart.east) -- (6.15,1.65);
 \node at (7.20,1.65) {$\cdots$};
 \draw[arrow,densely dashed] (8.25,1.65) -- (graphend.west);
 \draw[guide] (actualstart.south) -- (graphstart.north);
 \node[fill=white,align=center] at (7.20,2.23)
   {same event indices;\\unbounded return count};
 \node[anchor=west] at (0,.82)
   {Same actual base; different complete fibres.};
 \node[anchor=west] at (0,.22)
   {Same-index comparison: Lemma~\ref{sel:sharp-phase}.};
\end{tikzpicture}
\caption{From the physical preparation to the same-base return
comparison. (a) Both preparatory flights and the intervening signed
contact are generated by the original trajectory.
(b) The graph starts at the actual third-touchdown base, with a small
complete fibre discrepancy. The actual normal memory is retained;
the two starting states are not identical.
Arrows and spacing are schematic, not physical time.}
\label{fig:entry-comparison}
\end{figure}


\subsection{Normalized touchdown coordinates and the finite-scale reference}

The return sends an input scale $\delta$ to
$D\asymp\delta^{3/4}$. Under the bounds below, its tangential
correction has size $O(\delta)=O(D^{4/3})$. The normal correction
has size $O(\delta D)=O(D^{7/3})$ after the actual touchdown equation
eliminates the larger release-velocity terms. These cancellations,
proved in Appendix~\ref{sec:detuned-return}, motivate the powers
$4/3$ and $7/3$ in the normalized state.

Use the slips $u_r,u_0=u_r+r$ defined in Section~\ref{sec:model},
and put $r_0=c_\lambda/3$. At a touchdown with $p=-a\delta<0$, the independent coordinates are
\begin{equation}
 \begin{gathered}
 y=z=\eps\delta^2\zeta,\qquad
 u_r=d_\lambda\delta+d_\lambda\delta^{4/3}\eta_U^r,\qquad
 v=-r_0\delta^2-r_0\delta^{7/3}\eta_W,\\
 \xi^r=(\zeta,\eta_U^r,\eta_W),\qquad
 W=1+\delta^{1/3}\eta_W.
 \label{loc:tuple}
 \end{gathered}
\end{equation}
If $p(\phi_\delta)=-a\delta$ on the chosen positive-cosine branch,
the physical inverse is
\begin{equation}
 (\phi,y,\omega,v,z)
 =\left(\phi_\delta,\eps\delta^2\zeta,
       \frac{u_r+r-U_*}{\cos\phi_\delta},v,\eps\delta^2\zeta\right).
 \label{loc:inverse}
\end{equation}
Thus $b$ and $\eta_B=(-b-c_\lambda\delta)/(c_\lambda\delta^{4/3})$
are dependent coordinates. Fix a closed ball $\Kcal$ in the independent
coordinates with $0$ and all $\xi_{*,\lambda}$ in its interior.
Lemma~\ref{lem:td2-complete-preparation} bounds the original second-touchdown
independent tuples uniformly. Choose $M_{\rm in}$ once so that its ball
contains these tuples and $\Kcal$ strictly in its interior.
This choice does not change the fixed slip threshold $M_{\rm sl}$.
Fix $0<\gamma_{\rm cl}<\frac12\min_\lambda r_0$.
The physical input domain is
\begin{equation}
 |\xi^r|\le M_{\rm in},\quad r_0W-\zeta\ge\gamma_{\rm cl}>0,
 \quad 0<\eps\le C_0\delta^2,
 \quad 0\le r\le A\delta^{4/3},\qquad e=\eps\delta^{-3/2}.
 \label{loc:domain}
\end{equation}
The closing margin is $\dot z-v=\delta^2(r_0W-\zeta)$.
There is no positive lower bound on $e$ or $r$.
Taylor expansion in the chart $(p,y,u_0,v)$ gives
\begin{equation}
 \eta_B=\frac{g_b}{c_\lambda}
       \left(\eta_U^r+\frac{r}{d_\lambda\delta^{4/3}}\right)
       +O(\delta^{2/3}),\qquad g_b=-d_\lambda b_u(G),
 \label{loc:dependent}
\end{equation}
where $b_u$ is the derivative with respect to $u_0$.
The identity $-ab_p(G)+d_\lambda b_u(G)=-c_\lambda$ supplies the
leading cancellation; see Lemma~\ref{ret:inverse}.

The limiting return is defined at a fixed positive input scale
$\delta$, with $\eps,r\to0$. Physical displacement and memory then
vanish, but their normalized values still determine the release speed.
The later expansion $\delta\downarrow0$ is used to construct the graph
and its phase. To define this exact return $T_{0,\lambda}=(f_0,h_0)$,
set $r=0$ and use identical values of the independent coordinates $\xi$.
Put $\sigma=\sqrt\delta$ and
$h_\delta=\mu A(\phi_\delta)-B(\phi_\delta)$.
Use the compression clock $\tau_c=\sqrt\delta(t-t_{\rm in})/\eps$
and scalings $(y,z)=\eps\delta^{3/2}(Y,Z)$, $v=\delta^2V_c$
and $u_0=\delta\mathcal U$.
With primes denoting $\tau_c$ derivatives, the saturated compression
reference solves
\begin{equation}
 \begin{gathered}
 Y'=V_c,\quad \sigma Z'=Y-2Z,\quad
 V_c'=-a(Z-Y),\quad \mathcal U'=-h_\delta(Z-Y),\\
 (Y,Z,V_c,\mathcal U)(0)=
 (\sigma\zeta,\sigma\zeta,-r_0W,
            d_\lambda(1+\delta^{1/3}\eta_U)).
 \label{loc:compression}
 \end{gathered}
\end{equation}
At its first arrest $\mathcal U=0$, switch to the elapsed release
clock $\tau_r=(t-t_{\rm arrest})/\eps$ and velocity
$v=\delta^{3/2}V$. Thus $V=\sqrt\delta V_c$ at arrest, while
$Y$ and $Z$ keep their compression scalings. With primes now denoting
$\tau_r$ derivatives, the inherited normal motion $N=(Y,V,Z)$ solves
\begin{equation}
 N'=A_DN,\qquad
 A_D=\begin{pmatrix}0&1&0\\-1/2&0&1/2\\1&0&-2\end{pmatrix},
 \qquad N(0)=(Y,\sigma V_c,Z)_{\rm arrest}.
 \label{loc:normal-reference}
\end{equation}
If $R_0$ is its first force zero and $Q=V(R_0)$, reconstruct the
physical release as
$(p,y,u_0,v,z)=(-a\delta,0,0,\delta^{3/2}Q,0)$.
Its exact free mechanical flow to the first subsequent touchdown
defines $T_{0,\lambda}$. This retains the finite $\sqrt\delta$
compression and its inherited memory; the zero physical memory is
introduced only in the release reconstruction.

For the initial state $(-2,0,-1)$,
Lemma~\ref{con:normal-opening} gives $Y=Z=-\gamma_\#$ and
$V=\nu_\#$ at the first release, with
$\nu_\#>\gamma_\#>0$. Define
\begin{equation}
 \begin{gathered}
 q_\lambda=\nu_\#\sqrt{r_0d_\lambda/h+ad_\lambda^2/(2h^2)},
 \qquad K_\lambda=\sqrt{6q_\lambda/c_\lambda},\\
 \varkappa_W=\tfrac12+\frac{3g_b}{2c_\lambda},\qquad
 \xi_{*,\lambda}=K_\lambda^{-4/3}(0,-1,-\varkappa_W).
 \label{loc:constants}
 \end{gathered}
\end{equation}

\subsection{Preparation of the original trajectory}

Preparation serves two purposes. Its third-touchdown state must be
accurate enough for repeated return comparison and for the entrance
phase in Section~\ref{sec:graph}. The whole preceding trajectory
must also approach the prescribed incoming arc, so that later returns
can be attached to it. We first record the state estimates and then
the bounds on the intervening motion.

The capture and critical-passage estimates give a first release with
$v_r=\eps(q_{\rm in}+o(1))$, negative normal memory and a strict
opening margin. Here $q_{\rm in}>0$ depends only on $(\lambda,\theta)$,
while $p_r/\sqrt\eps\to H=ak/\gamma$ and
$b_r/\sqrt\eps\to-\ell H$
(Theorem~\ref{thm:detuned-release}). With $c=c_\lambda$, the
preparation constants are
\begin{equation}
 \begin{gathered}
 cT_f^2+3\ell HT_f=6q_{\rm in},\quad T_f>0,
 \qquad \pi=H-aT_f,\\
 S=d_\lambda T_f,\quad Q=T_f(\ell H/2+cT_f/3),\quad
 A_{\rm reb}^2=QS/h-\pi S^2/(2h^2),\\
 R=\nu_\#A_{\rm reb},\qquad C_R=\sqrt{6R/c}.
 \label{loc:preparation-constants}
 \end{gathered}
\end{equation}
The scale $r=k\sqrt\eps$ balances this first flight: with
$t=\sqrt\eps\,T$, the terms $v_rt$, $b_rt^2/2$ and $-ct^3/6$
in \eqref{prep:gap-expansion} all have order $\eps^{3/2}$.
The smaller memory value still determines opening on $t=O(\eps)$
through $q_{\rm in}+Y_r^0>0$. Thus the same original preparation
supplies the positive release coefficient and inherited memory for
both exceptional flights.
At the first touchdown, $p_1=p(t_1)$ satisfies
$p_1/\sqrt\eps\to\pi$. Here $\sqrt\eps$ is the preparation scale,
and $\pi$ may be positive, zero or negative. Its arrest margin is
\[
 2hQ-\pi S
 =T_f\{\gamma H+(2hc/3+ad_\lambda)T_f\}>0.
\]
The signed contact theorem applies with $r=k\sqrt\eps$ and yields
$v_L=\eps^{3/4}(R+o(1))$ and
$y_L=z_L=-\eps^{7/4}(\gamma_\#A_{\rm reb}+o(1))$.
In particular both $R$ and $R-\gamma_\#A_{\rm reb}$ have common
positive lower bounds.

The ensuing touchdown occurs after an elapsed time $T$ and has scale
$D=-p(T)/a=C_R\eps^{3/8}(1+o(1))$. To obtain its complete tuple,
one must eliminate the release velocity before dividing by the finer
normal scale. A separate release-velocity approximation would be
amplified by that division. At the actual root $y(T)=z(T)$, direct integration gives
\begin{equation}
 -v(T)-\frac c3D^2
 =\frac{y_L-z(T)}T-\frac12b_LT+\frac c3(T^2-D^2)
 -\frac1{2T}\int_0^T(T^2-t^2)(\dot b(t)+c)\,\dd t.
 \label{loc:touchdown-cancellation}
\end{equation}
The tangential and acceleration numerators are $O(\sqrt\eps)$,
and this normal numerator is $O(\eps^{7/8})$.
Their denominators $D^{4/3},D^{4/3},D^{7/3}$ have exactly these
orders. Thus the actual second touchdown has a bounded full tuple,
negative memory, a strict closing margin, and
$r/D^{4/3}=kC_R^{-4/3}(1+o(1))$; see
Lemma~\ref{lem:td2-complete-preparation}.

\begin{proposition}[The original prepared entrance]
\label{loc:entrance}
Under the standing model convention, all the following estimates are
uniform on the prescribed compact label classes.
The original trajectory has an admitted third touchdown with
\begin{equation}
 \begin{gathered}
 \delta_2=C_R\eps^{3/8}(1+o(1)),\qquad
 \delta_3\asymp\eps^{9/32},\quad
 r/\delta_3^{4/3}=O(\eps^{1/8}),\\
 \left|\log\frac{\delta_3}{K_\lambda\delta_2^{3/4}}\right|
 +|\xi_3^r-\xi_{*,\lambda}|+
 |\eta_{B,3}-(g_b/c)\eta_{U,*,\lambda}|
 \le C\eps^{3/32}.
 \label{loc:third-entry}
 \end{gathered}
\end{equation}
Here $\eta_{U,*,\lambda}$ is the slip component of $\xi_{*,\lambda}$.
The third touchdown retains the full cone, inverse branch, negative
normal memory, strict closing margin and all inequalities of
\eqref{loc:domain}, so it is a new admitted input. The preparation scale
$\sqrt\eps$ is followed by touchdown scales $\eps^{3/8}$ and
$\eps^{9/32}$; only at the last does the normalized state approach
$\xi_{*,\lambda}$.
All estimates are uniform on the fixed compact label classes and
are obtained from the original physical initial state.


The entrance constant
\begin{equation}
 \chi_\lambda(\mathfrak a)=\chi_{\lambda,\theta}(k)
   =3\log K_\lambda-\tfrac34\log C_R(\lambda,k,\theta),
 \qquad \mathfrak a=(k,\theta,\nu,\mathbf n),
 \label{sel:chi}
\end{equation}
is continuous, bounded and independent of $\nu,\mathbf n$. Moreover,
\begin{equation}
 \log(K_\lambda^4/\delta_3^{\mathfrak a})
 =\tfrac9{32}\log(1/\eps)+\chi_\lambda(\mathfrak a)
                      +r_{\mathfrak a}^{\log}(\eps),
 \qquad \sup_{\mathfrak a}|r_{\mathfrak a}^{\log}(\eps)|\to0.
 \label{loc:log-entry}
\end{equation}
On each fixed closed incoming rigid subarc the original motion
approaches that prescribed arc. For each fixed small $\eta>0$,
its remaining prefix from $p=\eta$ through $\TD_3$ is contained in
$G+B_{C\eta+o(1)}$ in the full physical state. The motion has positive
$p$ up to the descending-slip threshold, and its mechanical distance
from $G$ tends uniformly to zero from that threshold through $\TD_3$.
\end{proposition}
\begin{proof}
Lemma~\ref{lem:td2-complete-preparation} and
Corollary~\ref{cor:original-to-td2} give the complete second input
and its uniform scale asymptotic. Applying Theorems~\ref{thm:detuned-value-bridge}
and~\ref{thm:detuned-output-admission} to this input, as in
Corollary~\ref{cor:td3-admission}, gives \eqref{loc:third-entry}
and the stated admissibility of the returned state.
The normalized first-release data are continuous and have common
positive margins by Lemma~\ref{lem:critical-pullback} and
Theorem~\ref{thm:detuned-release}. The formulas
\eqref{loc:preparation-constants} therefore give common positive
bounds and continuity for $C_R$. Set
\[
 e_{\mathfrak a}(\eps)
 =\log\frac{\delta_2^{\mathfrak a}}
                 {C_R(\mathfrak a)\eps^{3/8}}.
\]
The second-touchdown root estimates are uniform, so
$\sup_{\mathfrak a}|e_{\mathfrak a}|\to0$.
Taking logarithms in \eqref{loc:third-entry} gives
$r_{\mathfrak a}^{\log}=-\tfrac34 e_{\mathfrak a}
+O(\eps^{3/32})$, proving \eqref{loc:log-entry}.
The second-touchdown remainder needs no specified rate.

On a closed incoming arc $p\ge\eta$, the unilateral capture estimate
in the proof of Lemma~\ref{lem:detuned-threshold} bounds the full scaled normal
state through any initial departures and recontacts. Capture takes
a common finite fast time, produces an $O(\eps)$ slow increment,
and is followed by an $O(\eps)$ integrated slow-field discrepancy
from the prescribed saturated rigid orbit; see
\eqref{ent:fixed-initialization}. Gronwall's inequality gives convergence
to that particular rigid subarc. Throughout it
$|y|+|z|=O_\eta(\eps^2)$ and $|v|=O_\eta(\eps)$, so the fast
excursion collapses to its original physical starting point.

For the remaining incoming segment, the entry estimates
\eqref{ent:safe-tracking} control the normal state and its effect
on the slow variables up to the descending-slip threshold.
They apply on every stopped interval before that threshold, including
the final $p\asymp\sqrt\eps$ annulus. The weak-coordinate bounds
and analytic chart put this segment in $G+B_{C\eta+o(1)}$; it remains
at positive $p$ until that threshold. The stopped comparison in
Theorem~\ref{thm:detuned-release} controls every prefix through its first
force zero, with mechanical distance
$O(\sqrt\eps+\eps(1+|\log\eps|))$ from $G$.
Proposition~\ref{prep:first-flight} gives duration and mechanical state
distance $O(\sqrt\eps)$ through the first exceptional touchdown.
The ensuing contact uses the first application of
Theorem~\ref{con:contact-release}, with scale $\sqrt\eps$ and
$r=O(\sqrt\eps)$; its complete prefix bound is $O(\sqrt\eps)$.
The second exceptional flight has duration and mechanical state
distance from $G$ both $O(\eps^{3/8})$. The full-prefix estimates
leading to Corollary~\ref{cor:original-to-td2} control these passages
and supply the actual complete $\TD_2$ cone.
Apply the contact prefix estimate \eqref{con:prefix} and
Lemma~\ref{lem:paired-free-normalization} to the
$\TD_2\to\TD_3$ return at $\delta_2\asymp\eps^{3/8}$.
Its mechanical distance from $G$ is $O(\eps^{9/32})$ and its memory
tends to zero. These bounds include both exceptional flights, their
contacts and the passage to $\TD_3$, and prove the stated original-prefix
localization.

\end{proof}

\subsection{The return estimate and its domains}

Let $\|\cdot\|_{1,*}$ include values and first derivatives with
respect to $\delta\partial_\delta$ and $\xi$.
The following conclusions collect
Proposition~\ref{prop:limiting-factorization} and
Theorems~\ref{thm:detuned-value-bridge}--\ref{thm:detuned-output-admission}.
Write the actual return as
\[
 T_{\eps,r}(\delta,\xi^r)=(D,\xi^{r,+})
   =\bigl(f_{\eps,r}(\delta,\xi^r),h_{\eps,r}^r(\delta,\xi^r)\bigr).
\]
The reference derivatives construct the invariant graph; the value
error controls the actual iterates; the output conditions allow the
next return. All constants are uniform on the fixed compact model
classes and the prescribed finite input bounds.

\begin{proposition}[The complete local return]
\label{loc:return}
For sufficiently small $\delta$, the following statements hold.
\begin{enumerate}[label=\textup{(\roman*)}]
\item \emph{Reference regularity.} On a bounded convex auxiliary fibre
domain $\Omega$ containing the fixed independent input ball and $\Kcal$,
the exact reference has a $C^1$ extension with
\begin{equation}
 \begin{aligned}
 \left\|\log\frac{f_0}{K_\lambda\delta^{3/4}}\right\|_{1,*}
       +\|h_0-\xi_{*,\lambda}\|_{1,*}&\le C\delta^{1/4},\\
 |D_\xi\log f_0|+|D_\xi h_0|&\le C\delta^{1/3}.
 \end{aligned}
 \label{loc:factorization}
\end{equation}
The reference map and its state derivatives are jointly continuous in
$\lambda$. The memory component of $h_0$ vanishes.
\item \emph{Compliant value approximation.} On the physical domain
\eqref{loc:domain},
the actual first contact--flight return exists, retains its negative
memory and transverse first-event margins, and satisfies
\begin{equation}
 \left|\log\frac{f_{\eps,r}(\delta,\xi^r)}{f_0(\delta,\xi^r)}\right|
 +|h_{\eps,r}^r(\delta,\xi^r)-h_0(\delta,\xi^r)|
 \le C(e+r/\delta).
 \label{loc:value}
\end{equation}
\item \emph{Admissibility of the returned state.} There are common radii
$0<\delta_{\rm reset}<\delta_{\rm g}<\delta_{\rm call}$ such that
a physical input satisfying \eqref{loc:domain} with
$\delta\le\delta_{\rm reset}$ has output $D$ satisfying
\begin{equation}
 \begin{gathered}
 \delta<D<\delta_{\rm g},\quad \xi^{r,+}\in\operatorname{int}\Kcal,
 \quad \eps<C_0D^2,\quad r\le AD^{4/3},\\
 \left|\log\frac D{K_\lambda\delta^{3/4}}\right|
       +|\xi^{r,+}-\xi_{*,\lambda}|
       \le C(\delta^{1/4}+e+r/\delta).
 \label{loc:admission}
 \end{gathered}
\end{equation}
It has $u_r^+>0$, $v^+<0$, $y^+=z^+<0$, the same inverse branch,
a bounded dependent coordinate and the strict closing margin.
Another return requires $D\le\delta_{\rm reset}$.
The outer radius may be reduced first, with the reset radius chosen
afterwards.
\end{enumerate}
\end{proposition}

The auxiliary extension in \eqref{loc:factorization} continues the
signed linear reference even when initial force is negative.
For physical inputs, the proof of \eqref{loc:value} first compares the
actual return with a reference $R^r$ reconstructed at slip
$u_{0,L}=r$. This costs $Ce$. The references $R^r$ and $T_0$ share
the entire finite-scale compression and inherited release, and their
free flights differ by $Cr/\delta$. Finally the actual output is
recoded by
\[
 h_{\eps,r}^r=h_{\eps,r}^0-
                  \frac{r}{d_\lambda D^{4/3}}(0,1,0).
\]
The paired free-flight estimate has already established
$D^{4/3}\asymp\delta$, so this last cost is also $Cr/\delta$.
The contact comparison retains both tail moments and the inherited
normal state. In the elapsed release time $\rho$ of
Appendix~\ref{sec:contact}, put $w=u_r/\eps$, $L=Z-Y$ and
$\vartheta=B/A$. The normal-velocity projection is
\[
 \begin{gathered}
 P=\widehat V-e\vartheta w,\qquad \widetilde N=(Y,P,Z),\\
 \widetilde N'=A_D\widetilde N+eF,\qquad
 F=(\vartheta w,b-\vartheta a_T-\vartheta'w,0).
 \end{gathered}
\]
where primes denote $\rho$ derivatives and
$\vartheta'=\eps\vartheta_\phi\dot\phi$.
The exact identity $C-B^2/A=1/2$ cancels the term containing $w'$
when the normal and slip equations are combined. Thus $F$ is bounded
on the stopped boxes without estimating the stiff slip derivative.
The full stopped-prefix estimates and first-release brackets are
proved in \eqref{con:projection}--\eqref{con:release-mechanics}.
Only the exact reference return is differentiated in
\eqref{loc:factorization}; \eqref{loc:value} compares values of the
actual compliant map.

